% =====================================================================
\newpage
\section{Особенности реализации}
\label{sec:impl}

Программа написана на языке Java и состоит из двух классов:
\begin{itemize}
	\item \texttt{CellularAutomaton} --- логика автомата: хранение поля, правило перехода, шаг эволюции, начальные условия, обнаружение цикла;
	\item \texttt{Main} --- консольный интерфейс: ввод параметров с проверкой, выбор начальных условий, вывод поля и итогов.
\end{itemize}
Полный исходный код приведен в \hyperref[app:code]{Приложении}.

\subsection{Класс CellularAutomaton}

Класс содержит следующие поля:
\begin{itemize}
	\item \texttt{long rule} --- номер правила;
	\item \texttt{int size} --- ширина квадратного поля;
	\item \texttt{boolean[][] cells} --- текущее поколение;
	\item \texttt{boolean[][] next} --- буфер для следующего поколения;
	\item \texttt{int generation} --- номер текущего поколения;
	\item \texttt{Map<BitSet, Integer> history} --- все встреченные конфигурации и номера поколений, в которых они впервые появились;
	\item \texttt{boolean cycleFound}, \texttt{int transientLength}, \texttt{int period} --- признак найденного цикла, предпериод и период.
\end{itemize}

\subsubsection*{Конструктор CellularAutomaton()}

Конструктор вычисляет номер правила по формуле из задания и создает два пустых поля заданного размера. Приведение первого множителя к \texttt{long} гарантирует, что все произведение вычисляется в 64-битной арифметике без переполнения. Код представлен в листинге~\ref{lst:ctor}.

Входные данные:
\begin{itemize}
	\item \texttt{int numberOption, constValue, year, month, day} --- номер варианта, константа 11 и дата рождения из формулы номера правила;
	\item \texttt{int size} --- ширина поля.
\end{itemize}

Выходные данные:
\begin{itemize}
	\item созданный объект автомата с вычисленным номером правила;
	\item пустое поле (все клетки мертвые, номер поколения 0).
\end{itemize}

\begin{lstlisting}[caption={Конструктор класса CellularAutomaton}, label={lst:ctor}]
public CellularAutomaton(int numberOption, int constValue, int year, int month, int day, int size) {
	this.rule = (long) numberOption * constValue * year * month * day;
	this.size = size;
	this.cells = new boolean[this.size][this.size];
	this.next = new boolean[this.size][this.size];
}
\end{lstlisting}

\subsubsection*{Метод bit()}

Метод \texttt{bit} возвращает $k$-й бит числа $x$: число сдвигается вправо на $k$ разрядов, после чего побитовым «И» с единицей выделяется младший бит. Код представлен в листинге~\ref{lst:bit}.

Входные данные:
\begin{itemize}
	\item \texttt{long x} --- число;
	\item \texttt{int k} --- номер бита.
\end{itemize}

Выходные данные:
\begin{itemize}
	\item значение бита, 0 или 1.
\end{itemize}

\begin{lstlisting}[caption={Метод bit()}, label={lst:bit}]
static int bit(long x, int k) {
	long shifted = x >> k;
	long lastBit = shifted & 1;
	return (int) lastBit;
}
\end{lstlisting}

\subsubsection*{Метод newState()}

Метод \texttt{newState} применяет \texttt{bit} к номеру правила, то есть вычисляет значение функции перехода $f$ на наборе с номером \texttt{idx}. Код представлен в листинге~\ref{lst:newstate}.

Входные данные:
\begin{itemize}
	\item \texttt{int idx} --- индекс окрестности ($0..31$).
\end{itemize}

Выходные данные:
\begin{itemize}
	\item новое состояние клетки, 0 или 1.
\end{itemize}

\begin{lstlisting}[caption={Метод newState()}, label={lst:newstate}]
int newState(int idx) {
	return bit(rule, idx);
}
\end{lstlisting}

\subsubsection*{Метод wrap()}

Метод \texttt{wrap} реализует тороидальные граничные условия: переводит любую координату в диапазон $0..size-1$. Код представлен в листинге~\ref{lst:wrap}.

Входные данные:
\begin{itemize}
	\item \texttt{int coord} --- координата (строка или столбец), возможно, вышедшая за край поля.
\end{itemize}

Выходные данные:
\begin{itemize}
	\item координата \texttt{coord mod size} в диапазоне $0..size-1$.
\end{itemize}

\begin{lstlisting}[caption={Метод wrap()}, label={lst:wrap}]
private int wrap(int coord) {
	return (coord % size + size) % size;
}
\end{lstlisting}

\subsubsection*{Метод getWrapped()}

Метод \texttt{getWrapped} возвращает состояние клетки, приводя обе координаты методом \texttt{wrap}. Код представлен в листинге~\ref{lst:getwrapped}.

Входные данные:
\begin{itemize}
	\item \texttt{int row, col} --- координаты клетки, возможно, вышедшие за край поля.
\end{itemize}

Выходные данные:
\begin{itemize}
	\item состояние клетки \texttt{cells[wrap(row)][wrap(col)]} (\texttt{true} --- живая).
\end{itemize}

\begin{lstlisting}[caption={Метод getWrapped()}, label={lst:getwrapped}]
private boolean getWrapped(int row, int col) {
	return cells[wrap(row)][wrap(col)];
}
\end{lstlisting}

\subsubsection*{Метод neighborhoodIndex()}

Метод собирает индекс окрестности фон Неймана клетки $(row, col)$ в порядке $C\,N\,E\,W\,S$. Состояние каждой из пяти клеток читается с учетом тора и переводится в 0 или 1 методом \texttt{toBit}, после чего биты складываются с весами 16, 8, 4, 2, 1. Код представлен в листинге~\ref{lst:index}.

Входные данные:
\begin{itemize}
	\item \texttt{int row, col} --- координаты клетки.
\end{itemize}

Выходные данные:
\begin{itemize}
	\item индекс окрестности $idx = 16C + 8N + 4E + 2W + S$ в диапазоне $0..31$.
\end{itemize}

\begin{lstlisting}[caption={Метод neighborhoodIndex()}, label={lst:index}]
int neighborhoodIndex(int row, int col) {
	int n = toBit(getWrapped(row - 1, col));
	int c = toBit(getWrapped(row, col));
	int s = toBit(getWrapped(row + 1, col));
	int w = toBit(getWrapped(row, col - 1));
	int e = toBit(getWrapped(row, col + 1));
	return c * 16 + n * 8 + e * 4 + w * 2 + s;
}
\end{lstlisting}

\subsubsection*{Метод toBit()}

Вспомогательный метод \texttt{toBit} переводит \texttt{boolean} в число: \texttt{true} $\to$ 1, \texttt{false} $\to$ 0. Код представлен в листинге~\ref{lst:tobit}.

Входные данные:
\begin{itemize}
	\item \texttt{boolean value} --- состояние клетки.
\end{itemize}

Выходные данные:
\begin{itemize}
	\item 1 или 0.
\end{itemize}

\begin{lstlisting}[caption={Метод toBit()}, label={lst:tobit}]
private static int toBit(boolean value) {
	return value ? 1 : 0;
}
\end{lstlisting}

\subsubsection*{Метод step()}

Метод выполняет один шаг эволюции. Для каждой клетки вычисляется индекс окрестности по текущему полю \texttt{cells}, новое состояние записывается в буфер \texttt{next}. После обработки всех клеток ссылки на массивы меняются местами: новое поколение становится текущим, а старый массив используется как буфер на следующем шаге. Копирования массивов при этом не происходит. Перед первым шагом в историю записывается начальная конфигурация, а после каждого шага выполняется проверка на повторение конфигурации. Сложность одного шага --- $O(n^2)$. Код представлен в листинге~\ref{lst:step}.

Входные данные:
\begin{itemize}
	\item используется текущее поле \texttt{cells}.
\end{itemize}

Выходные данные:
\begin{itemize}
	\item поле \texttt{cells} заменяется следующим поколением;
	\item номер поколения \texttt{generation} увеличивается на 1;
	\item при повторении конфигурации заполняются поля:
	\begin{itemize}
		\item \texttt{cycleFound} --- признак найденного цикла;
		\item \texttt{transientLength} --- предпериод;
		\item \texttt{period} --- период.
	\end{itemize}
\end{itemize}

\begin{lstlisting}[caption={Метод step()}, label={lst:step}]
public void step() {
	if (history.isEmpty()) {
		history.put(snapshot(), generation);
	}

	for (int i = 0; i < size; i++) {
		for (int j = 0; j < size; j++) {
			int idx = neighborhoodIndex(i, j);
			next[i][j] = newState(idx) == 1;
		}
	}

	boolean[][] tmp = cells;
	cells = next;
	next = tmp;

	generation++;
	checkCycle();
}
\end{lstlisting}

\subsubsection*{Метод snapshot()}

Метод \texttt{snapshot} упаковывает поле в \texttt{BitSet}: клетка $(row, col)$ соответствует биту $row \cdot size + col$, т.\,е. на клетку тратится один бит. Методы \texttt{equals} и \texttt{hashCode} класса \texttt{BitSet} сравнивают содержимое, поэтому его можно использовать как ключ хеш-таблицы, в отличие от массива \texttt{boolean[][]}. Код представлен в листинге~\ref{lst:snapshot}.

Входные данные:
\begin{itemize}
	\item используется текущее поле \texttt{cells}.
\end{itemize}

Выходные данные:
\begin{itemize}
	\item \texttt{BitSet} из $size^2$ битов, в котором установлены биты живых клеток.
\end{itemize}

\begin{lstlisting}[caption={Метод snapshot()}, label={lst:snapshot}]
private BitSet snapshot() {
	BitSet bits = new BitSet(size * size);
	for (int row = 0; row < size; row++) {
		for (int col = 0; col < size; col++) {
			if (cells[row][col]) {
				bits.set(row * size + col);
			}
		}
	}
	return bits;
}
\end{lstlisting}

\subsubsection*{Метод checkCycle()}

Для обнаружения цикла все встреченные конфигурации хранятся в хеш-таблице \texttt{history} (конфигурация $\to$ номер поколения). Метод \texttt{checkCycle} пытается добавить текущую конфигурацию в таблицу методом \texttt{putIfAbsent}. Если такая конфигурация уже встречалась в поколении $prev$, метод возвращает $prev$, и тогда предпериод равен $prev$, а период --- $generation - prev$. После обнаружения цикла конфигурации больше не сохраняются. Сравнение конфигураций точное (без риска коллизий), а затраты памяти составляют $O(T \cdot n^2 / 8)$ байт, где $T$ --- число выполненных шагов. Код представлен в листинге~\ref{lst:cycle}.

Входные данные:
\begin{itemize}
	\item используются текущее поле, номер поколения \texttt{generation} и таблица \texttt{history}.
\end{itemize}

Выходные данные:
\begin{itemize}
	\item текущая конфигурация добавляется в \texttt{history};
	\item предпериод \texttt{transientLength};
	\item период \texttt{period}.
\end{itemize}

\begin{lstlisting}[caption={Метод checkCycle()}, label={lst:cycle}]
private void checkCycle() {
	if (cycleFound) {
		return;
	}
	Integer prev = history.putIfAbsent(snapshot(), generation);
	if (prev != null) {
		cycleFound = true;
		transientLength = prev;
		period = generation - prev;
	}
}
\end{lstlisting}

\subsubsection*{Метод randomize()}

Метод заполняет поле случайным образом: каждая клетка становится живой с вероятностью \texttt{density}, независимо от остальных. Для этого используется равномерно распределенное число из $[0, 1)$, возвращаемое \texttt{Random.nextDouble()}: оно меньше \texttt{density} ровно с вероятностью \texttt{density}. Ожидаемое число живых клеток равно $density \cdot n^2$. Генератор псевдослучайных чисел инициализируется заданным значением \texttt{seed}, поэтому одинаковый \texttt{seed} всегда дает одинаковое поле, что позволяет воспроизводить эксперименты. Код представлен в листинге~\ref{lst:random}.

Входные данные:
\begin{itemize}
	\item \texttt{double density} --- плотность живых клеток ($0..1$);
	\item \texttt{long seed} --- начальное значение генератора.
\end{itemize}

Выходные данные:
\begin{itemize}
	\item поле \texttt{cells} заполняется случайно.
\end{itemize}

\begin{lstlisting}[caption={Метод randomize()}, label={lst:random}]
public void randomize(double density, long seed) {
	Random random = new Random(seed);
	for (int row = 0; row < size; row++) {
		for (int col = 0; col < size; col++) {
			cells[row][col] = random.nextDouble() < density;
		}
	}
}
\end{lstlisting}

\subsubsection*{Метод getCell()}

Метод возвращает состояние клетки. Код представлен в листинге~\ref{lst:getcell}.

Входные данные:
\begin{itemize}
	\item \texttt{int row, col} --- координаты клетки в диапазоне $0..size-1$.
\end{itemize}

Выходные данные:
\begin{itemize}
	\item \texttt{true}, если клетка живая.
\end{itemize}

\begin{lstlisting}[caption={Метод getCell()}, label={lst:getcell}]
public boolean getCell(int row, int col) {
	return cells[row][col];
}
\end{lstlisting}

\subsubsection*{Метод setCell()}

Метод задает состояние клетки, используется при ручном вводе и при размещении клетки в центре поля. Код представлен в листинге~\ref{lst:setcell}.

Входные данные:
\begin{itemize}
	\item \texttt{int row, col} --- координаты клетки;
	\item \texttt{boolean value} --- новое состояние.
\end{itemize}

Выходные данные:
\begin{itemize}
	\item изменяется клетка поля.
\end{itemize}

\begin{lstlisting}[caption={Метод setCell()}, label={lst:setcell}]
public void setCell(int row, int col, boolean value) {
	cells[row][col] = value;
}
\end{lstlisting}

\subsubsection*{Метод clear()}

Метод делает все клетки поля мертвыми. Код представлен в листинге~\ref{lst:clear}.

Входные данные:
\begin{itemize}
	\item текущее поле \texttt{cells}.
\end{itemize}

Выходные данные:
\begin{itemize}
	\item поле очищено.
\end{itemize}

\begin{lstlisting}[caption={Метод clear()}, label={lst:clear}]
public void clear() {
	for (int row = 0; row < size; row++) {
		for (int col = 0; col < size; col++) {
			cells[row][col] = false;
		}
	}
}
\end{lstlisting}

\subsubsection*{Метод population()}

Метод подсчитывает число живых клеток поля. Код представлен в листинге~\ref{lst:population}.

Входные данные:
\begin{itemize}
	\item текущее поле \texttt{cells}.
\end{itemize}

Выходные данные:
\begin{itemize}
	\item число живых клеток от 0 до $size^2$.
\end{itemize}

\begin{lstlisting}[caption={Метод population()}, label={lst:population}]
public int population() {
	int count = 0;
	for (int row = 0; row < size; row++) {
		for (int col = 0; col < size; col++) {
			if (cells[row][col]) {
				count++;
			}
		}
	}
	return count;
}
\end{lstlisting}

\subsubsection*{Методы доступа}

Методы доступа возвращают значения закрытых полей класса (листинги~\ref{lst:getsize}--\ref{lst:getperiod}):
\begin{itemize}
	\item \texttt{getSize} --- ширина поля \texttt{size};
	\item \texttt{getGeneration} --- номер текущего поколения \texttt{generation};
	\item \texttt{isCycleFound} --- признак найденного цикла \texttt{cycleFound};
	\item \texttt{getTransientLength} --- предпериод \texttt{transientLength} (0, если цикл не найден);
	\item \texttt{getPeriod} --- период \texttt{period} (0, если цикл не найден).
\end{itemize}

\begin{lstlisting}[caption={Метод getSize()}, label={lst:getsize}]
public int getSize() {
	return size;
}
\end{lstlisting}

\begin{lstlisting}[caption={Метод getGeneration()}, label={lst:getgeneration}]
public int getGeneration() {
	return generation;
}
\end{lstlisting}

\begin{lstlisting}[caption={Метод isCycleFound()}, label={lst:iscyclefound}]
public boolean isCycleFound() {
	return cycleFound;
}
\end{lstlisting}

\begin{lstlisting}[caption={Метод getTransientLength()}, label={lst:gettransient}]
public int getTransientLength() {
	return transientLength;
}
\end{lstlisting}

\begin{lstlisting}[caption={Метод getPeriod()}, label={lst:getperiod}]
public int getPeriod() {
	return period;
}
\end{lstlisting}

\subsection{Класс Main}

Класс \texttt{Main} реализует консольный интерфейс программы.

\subsubsection*{Метод main()}

Метод запрашивает у пользователя ширину поля ($3..100$), количество итераций ($1..100\,000$), начальные условия, шаг вывода $k$ (печатается каждое $k$-е поколение) и необходимость остановки при обнаружении цикла. Затем выполняется заданное число итераций; печатаются начальное поле, каждое $k$-е поколение, последнее поколение и поколение, на котором обнаружен цикл. В конце выводится итог работы. Код представлен в листинге~\ref{lst:main}.

Входные данные:
\begin{itemize}
	\item значения вводятся пользователем.
\end{itemize}

Выходные данные:
\begin{itemize}
	\item поколения и итог работы выводятся в консоль.
\end{itemize}

Минимальная ширина поля равна 3, так как при $n \le 2$ на торе северный и южный (западный и восточный) соседи клетки совпадают, и окрестность фон Неймана вырождается.

\begin{lstlisting}[caption={Основной цикл метода main()}, label={lst:main}]
printGeneration(automaton);
for (int i = 1; i <= iterations; i++) {
	boolean hadCycle = automaton.isCycleFound();
	automaton.step();
	boolean cycleJustFound = !hadCycle && automaton.isCycleFound();

	if (i % printEvery == 0 || i == iterations || cycleJustFound) {
		printGeneration(automaton);
	}
	if (cycleJustFound && stopOnCycle) {
		break;
	}
}

printSummary(automaton);
\end{lstlisting}

\subsubsection*{Метод readInt()}

Метод \texttt{readInt} запрашивает целое число до тех пор, пока не будет введено корректное значение из диапазона $[min, max]$. При пустом вводе возвращается значение по умолчанию, которое показывается в подсказке в квадратных скобках. Ввод читается целой строкой и преобразуется в число методом \texttt{Integer.parseInt}; некорректный ввод перехватывается через исключение \texttt{NumberFormatException}. Код представлен в листинге~\ref{lst:readint}.

Входные данные:
\begin{itemize}
	\item \texttt{String prompt} --- подсказка;
	\item \texttt{int min, max} --- границы диапазона;
	\item \texttt{int def} --- значение по умолчанию.
\end{itemize}

Выходные данные:
\begin{itemize}
	\item введенное число из $[min, max]$ или \texttt{def}.
\end{itemize}

\begin{lstlisting}[caption={Метод readInt()}, label={lst:readint}]
static int readInt(String prompt, int min, int max, int def) {
	while (true) {
		String s = readLine(prompt + " [" + def + "]: ");
		if (s.isEmpty()) {
			return def;
		}
		try {
			int value = Integer.parseInt(s);
			if (value >= min && value <= max) {
				return value;
			}
		} catch (NumberFormatException ignored) {
		}
		println("  Введите целое число от " + min + " до " + max + ".");
	}
}
\end{lstlisting}

\subsubsection*{Метод readLine()}

Метод \texttt{readLine} выводит подсказку, читает строку и удаляет пробелы по краям. При окончании ввода (EOF) метод \texttt{readln} возвращает \texttt{null}, и тогда возвращается пустая строка. Код представлен в листинге~\ref{lst:readline}.

Входные данные:
\begin{itemize}
	\item \texttt{String prompt} --- подсказка.
\end{itemize}

Выходные данные:
\begin{itemize}
	\item введенная строка без пробелов по краям или пустая строка.
\end{itemize}

\begin{lstlisting}[caption={Метод readLine()}, label={lst:readline}]
static String readLine(String prompt) {
	String s = readln(prompt);
	return s == null ? "" : s.trim();
}
\end{lstlisting}

\subsubsection*{Метод initField()}

Метод выводит меню начальных условий и заполняет поле выбранным способом (листинг~\ref{lst:init}):
\begin{enumerate}
	\item случайно --- запрашиваются плотность живых клеток в процентах и \texttt{seed}, после чего вызывается метод \texttt{randomize};
	\item вручную --- метод \texttt{readManual};
	\item одна живая клетка в центре поля, клетка $(\lfloor size/2 \rfloor, \lfloor size/2 \rfloor)$;
	\item пустое поле.
\end{enumerate}

Входные данные:
\begin{itemize}
	\item \texttt{CellularAutomaton automaton} --- автомат, поле которого заполняется; номер способа вводится пользователем.
\end{itemize}

Выходные данные:
\begin{itemize}
	\item поле автомата заполнено выбранным способом.
\end{itemize}

\begin{lstlisting}[caption={Метод initField()}, label={lst:init}]
static void initField(CellularAutomaton automaton) {
	println();
	println("Начальные условия:");
	println("  1 - случайно");
	println("  2 - вручную");
	println("  3 - одна клетка в центре");
	println("  4 - пустое поле");
	int mode = readInt("Выбор", 1, 4, 1);

	switch (mode) {
		case 1 -> {
			int percent = readInt("Плотность живых клеток, % (0...100)", 0, 100, 50);
			long seed = readSeed();
			automaton.randomize(percent / 100.0, seed);
		}
		case 2 -> readManual(automaton);
		case 3 -> {
			int center = automaton.getSize() / 2;
			automaton.setCell(center, center, true);
		}
		default -> automaton.clear();
	}
}
\end{lstlisting}

\subsubsection*{Метод readSeed()}

Метод запрашивает начальное значение генератора случайных чисел до тех пор, пока не будет введено целое число или пустая строка. Если \texttt{seed} не введен, он берется из текущего времени (\texttt{System.nanoTime()}) и печатается, чтобы эксперимент можно было повторить. Код представлен в листинге~\ref{lst:seed}.

Входные данные:
\begin{itemize}
	\item значение вводится пользователем.
\end{itemize}

Выходные данные:
\begin{itemize}
	\item \texttt{long} --- введенный или сгенерированный \texttt{seed}.
\end{itemize}

\begin{lstlisting}[caption={Метод readSeed()}, label={lst:seed}]
static long readSeed() {
	while (true) {
		String s = readLine("Seed (пусто = случайный): ");
		if (s.isEmpty()) {
			long seed = System.nanoTime();
			println("Используется seed = " + seed);
			return seed;
		}
		try {
			return Long.parseLong(s);
		} catch (NumberFormatException e) {
			println("  Введите целое число или оставьте пустым.");
		}
	}
}
\end{lstlisting}

\subsubsection*{Метод readManual()}

Метод реализует ручной ввод начальной конфигурации: пользователь вводит строки поля сверху вниз, символы \texttt{+} и \texttt{1} обозначают живую клетку, любые другие --- мертвую. Строка короче ширины поля дополняется мертвыми клетками, лишние символы игнорируются, пустая строка завершает ввод (оставшиеся строки --- мертвые). Код представлен в листинге~\ref{lst:manual}.

Входные данные:
\begin{itemize}
	\item \texttt{CellularAutomaton automaton} --- автомат, поле которого заполняется; строки поля вводятся пользователем.
\end{itemize}

Выходные данные:
\begin{itemize}
	\item поле автомата заполнено по вводу пользователя.
\end{itemize}

\begin{lstlisting}[caption={Метод readManual()}, label={lst:manual}]
static void readManual(CellularAutomaton automaton) {
	int size = automaton.getSize();
	println("Вводите строки поля сверху вниз (до " + size + " строк).");
	println("Живая клетка: + или 1, мёртвая: . или 0. Пустая строка — закончить ввод.");
	automaton.clear();
	for (int row = 0; row < size; row++) {
		String line = readln(String.format("%3d> ", row));
		if (line == null || line.isBlank()) {
			break;
		}
		for (int col = 0; col < Math.min(line.length(), size); col++) {
			char ch = line.charAt(col);
			automaton.setCell(row, col, ch == '+' || ch == '1');
		}
	}
}
\end{lstlisting}

\subsubsection*{Метод printField()}

Метод печатает поле построчно: \texttt{+} --- живая клетка, пробел --- мертвая. Код представлен в листинге~\ref{lst:printfield}.

Входные данные:
\begin{itemize}
	\item \texttt{CellularAutomaton automaton} --- автомат.
\end{itemize}

Выходные данные:
\begin{itemize}
	\item поле выводится в консоль.
\end{itemize}

\begin{lstlisting}[caption={Метод printField()}, label={lst:printfield}]
static void printField(CellularAutomaton automaton) {
	for (int row = 0; row < automaton.getSize(); row++) {
		for (int col = 0; col < automaton.getSize(); col++) {
			print(automaton.getCell(row, col) ? "+ " : "  ");
		}
		println();
	}
	println();
}
\end{lstlisting}

\subsubsection*{Метод printGeneration()}

Метод печатает заголовок с номером поколения, числом живых клеток и, если цикл уже найден, предпериодом и периодом, а затем само поле методом \texttt{printField}. Код представлен в листинге~\ref{lst:printgen}.

Входные данные:
\begin{itemize}
	\item \texttt{CellularAutomaton automaton} --- автомат.
\end{itemize}

Выходные данные:
\begin{itemize}
	\item заголовок и поле выводятся в консоль.
\end{itemize}

\begin{lstlisting}[caption={Метод printGeneration()}, label={lst:printgen}]
static void printGeneration(CellularAutomaton automaton) {
	String header = "Поколение " + automaton.getGeneration() + " | живых: " + automaton.population();
	if (automaton.isCycleFound()) {
		header += " | цикл: предпериод = " + automaton.getTransientLength() + ", период = " + automaton.getPeriod();
	}
	println(header);
	printField(automaton);
}
\end{lstlisting}

\subsubsection*{Метод printSummary()}

Метод выводит итог: число выполненных итераций, число живых клеток, а также поколение, на котором найден цикл, предпериод и период либо сообщение о том, что цикл не обнаружен. Код представлен в листинге~\ref{lst:printsummary}.

Входные данные:
\begin{itemize}
	\item \texttt{CellularAutomaton automaton} --- автомат.
\end{itemize}

Выходные данные:
\begin{itemize}
	\item итог выводится в консоль.
\end{itemize}

\begin{lstlisting}[caption={Метод printSummary()}, label={lst:printsummary}]
static void printSummary(CellularAutomaton automaton) {
	int size = automaton.getSize();
	println("Итог");
	println("Выполнено итераций: " + automaton.getGeneration());
	println("Живых клеток: " + automaton.population() + " из " + size * size);
	if (automaton.isCycleFound()) {
		int t = automaton.getTransientLength();
		int p = automaton.getPeriod();
		println("Цикл найден на поколении " + (t + p));
		println("Предпериод: " + t);
		println("Период: " + p + (p == 1 ? " (неподвижная точка)" : ""));
	} else {
		println("Цикл за " + automaton.getGeneration() + " итераций не обнаружен.");
	}
}
\end{lstlisting}
